\section*{Problem 1}

设线性变换 $\mathcal{A} : \mathbb{R}^3 \to \mathbb{R}^3$ 的矩阵为
\[
A =
\begin{bmatrix}
2 & 1 & 0 \\
0 & 2 & 0 \\
0 & 0 & 3
\end{bmatrix}.
\]
判断下列子空间是否为 $\mathcal{A}$ 的不变子空间，并说明理由：
\[
W_1 = \Span\{(1,0,0)^\top\},\quad
W_2 = \Span\{(1,0,0)^\top,(0,1,0)^\top\},\quad
W_3 = \Span\{(1,1,0)^\top\}.
\]
\begin{proof}
只需证明对子空间的一组基向量，线性变换后的向量仍在该子空间内即可。

对 $W_1$，有
\[
A(1,0,0)^\top = (2,0,0)^\top \in W_1.
\]
对 $W_2$，有
\[
\begin{cases}
A(1,0,0)^\top = (2,0,0)^\top \in W_2, \\
A(0,1,0)^\top = (1,2,0)^\top \in W_2.
\end{cases}
\]
对 $W_3$，有
\[
A(1,1,0)^\top = (3,2,0)^\top \notin W_3.
\]
因此 \fbox{$W_1$ 和 $W_2$ 是 $\mathcal{A}$ 的不变子空间，而 $W_3$ 不是}。

\end{proof}

\section*{Problem 2}

设
\[
A =
\begin{bmatrix}
A_1 & A_2 \\
O & A_3
\end{bmatrix}
\]
为 $n$ 阶分块矩阵，其中 $A_1$ 为 $r$ 阶方阵。证明子空间
\[
W_1 = \left\{\left(\alpha_1,\ldots,\alpha_r,0,\ldots,0\right)^\top : \alpha_i \in \mathbb{R}\right\}
\]
是 $A$ 的不变子空间；并证明
\[
W_2 = \left\{\left(0,\ldots,0,\beta_{r+1},\ldots,\beta_n\right)^\top : \beta_j \in \mathbb{R}\right\}
\]
是 $A$ 的不变子空间当且仅当 $A_2 = O$。

\begin{proof}
不难写出 $W_1$ 的一组基向量为
\[
\left\{
\begin{pmatrix}\bvec e_{r,1}\\ O_{n-r}\end{pmatrix},
\begin{pmatrix}\bvec e_{r,2}\\ O_{n-r}\end{pmatrix},
\ldots,
\begin{pmatrix}\bvec e_{r,r}\\ O_{n-r}\end{pmatrix}
\right\}.
\]
对任意基向量，有
\[
A \begin{pmatrix}\bvec{e}_{r, i}\\ O_{n - r}\end{pmatrix} = \begin{bmatrix}
A_1 & A_2 \\
O_{\left(n - r\right) \times r} & A_3
\end{bmatrix}
\begin{pmatrix}
\bvec{e}_{r, i} \\
O_{n - r}
\end{pmatrix} = \begin{pmatrix}
A_1 \bvec{e}_{r, i} + A_2 O_{n - r} \\
O_{\left(n - r\right) \times r}\bvec e_{r,i} + A_3 O_{n - r}
\end{pmatrix} = \begin{pmatrix}
A_1 \bvec{e}_{r, i} \\
O_{n - r}
\end{pmatrix} \in W_1.
\]
因此 \fbox{$W_1$ 是 $A$ 的不变子空间}。

对 $W_2$，有
\[A \begin{pmatrix}O_r\\ \bvec{e}_{n-r,j}\end{pmatrix} = \begin{bmatrix}
A_1 & A_2 \\
O_{\left(n - r\right) \times r} & A_3
\end{bmatrix}
\begin{pmatrix}O_r \\
\bvec{e}_{n - r, j}
\end{pmatrix} = \begin{pmatrix}A_2 \bvec{e}_{n - r, j} \\
A_3 \bvec{e}_{n - r, j}
\end{pmatrix}.
\]
因此 $W_2$ 不变当且仅当 $A_2\bvec e_{n-r,j}=O_r$ 对每个 $j$ 都成立，也即 \fbox{$W_2$ 是 $A$ 的不变子空间当且仅当 $A_2 = O$}。

\end{proof}
